\documentclass[10pt,a4paper]{amsart}
\usepackage[margin=19mm]{geometry}
\usepackage{amsmath,amssymb,mathtools}
\usepackage[scr=rsfs,cal=euler]{mathalfa}
\usepackage{xfrac}
\usepackage[unicode,hidelinks,bookmarks=false]{hyperref}
\newcommand{\ie}{i.e.\ }
\newcommand{\cf}{cf.\ }
\newcommand{\resp}{respectively\ }
\newcommand{\Subsection}{\subsection}
\providecommand{\nobreakdash}{\nobreak\hbox{-}}
\setlength{\emergencystretch}{2em}
\multlinegap0pt
\usepackage{bm}
\usepackage{bbm}
\usepackage{dsfont}
\usepackage{xspace}
\usepackage{cancel}
\usepackage{mathtools}
\usepackage{amsthm}
\makeatletter
\@ifundefined{theorem}{\newtheorem{theorem}{Theorem}}{}
\makeatother
\title[The Efron-Stein inequality for identically distributed pairs]{The Efron-Stein inequality for identically distributed pairs}
\author{Jnaneshwar Baslingker, B\'alint Vir\'ag}
\date{Source: arXiv:2605.05388v1 (CC BY 4.0)}
\begin{document}
\maketitle

\noindent\emph{Source and licence.} The Efron-Stein inequality for identically distributed pairs. Version arXiv:2605.05388v1 (CC BY 4.0), distributed under the Creative Commons Attribution 4.0 International licence, as recorded in the arXiv licence metadata for this exact version and shown on the abstract page https://arxiv.org/abs/2605.05388v1. Full rights notice: licenses/arxiv-2605.05388-CC-BY-4.0.txt.\par\smallskip

\noindent\textit{Editorial scope.} Counted unit: the Theorem whose source label is \texttt{thm: finite support} in the article (source0.tex lines 160--165, source label \texttt{thm: finite support}), delivered together with the article's own complete original proof of it (source0.tex lines 166--198, one \texttt{proof} environment of 33 lines). No mathematics is written, repaired or re-derived: statement and proof are the article's own text line for line. Required context retained uncounted: none. Every definition, display and label of those lines is retained unchanged. Omitted from this package and not redistributed: 1--159, 199--228. Nothing inside a retained excerpt is omitted or altered: every retained body is delivered exactly as the article's lines are written, so no omission receipt of any kind accompanies this package. Citations: the retained excerpts carry no citation command at all, so no bibliography entry of the article is transcribed and no reference is redistributed. Reference adaptation: none was needed and none was made; every reference inside the delivered unit resolves inside it, so no unresolved reference or citation is produced. Printed slips kept literal: no printed word, symbol, hypothesis or number of the counted statement or of its proof was altered. Layout and class: the article's own class is \texttt{article} with packages including geometry, amsmath, amsthm, amssymb, natbib; the wrapper is the lane's pinned \texttt{amsart} editorial wrapper at 10pt on A4, which restates the theorem-like environments the retained text needs and adds the article's own macro definitions that the retained text needs (none), with extra wrapper packages (bm, bbm, dsfont, xspace, cancel, mathtools). The delivered numbering is the unit's own. Assets: no retained span contains a figure environment, a graphics-inclusion command, a PDF-page inclusion, a table or a TikZ picture; no image file is copied into this package, the package declares no asset dependency and no image file is redistributed with it. Licence: version arXiv:2605.05388v1 of this article is distributed under the Creative Commons Attribution 4.0 International licence, as recorded in the arXiv licence metadata for this exact version and shown on the abstract page https://arxiv.org/abs/2605.05388v1. A full rights notice is delivered at licenses/arxiv-2605.05388-CC-BY-4.0.txt. Characters: every retained excerpt is pure ASCII, so the delivered engine, XeLaTeX with its default Latin Modern text font, reports no missing character.
\begin{theorem}\label{thm: finite support}
For each \(i\) let \((X_i,Y_i)\) be an identically distributed pair, so that the pairs \((X_i,Y_i)\) are independent as \(i\) varies, and \(X_i\) can take at most \(k_i\) values. Let \(f\) be a bounded measurable function, and let \(Z_i=f(X_1,\ldots X_{i},Y_{i+1},\ldots Y_n)\). Then
$$
E(Z_n-Z_0)^2\le \rho(k) \sum_{i=1}^n E(Z_i-Z_{i-1})^2, \qquad \rho (k)\le \max_i k_i/2.
$$
\end{theorem}

\begin{proof}
It helps to allow complex-valued \(f\), in which case the result applies to the squared modulus:
$$
E|Z_n-Z_0|^2\le \frac{\kappa}{2}\sum_{j=1}^n E|Z_j-Z_{j-1}|^2
$$
First we consider the special case when all \(X_j\) is uniform on \(\mathbb Z_{k_j}=\{0,\ldots,k_j-1\}\), and \(Y_j=X_{j}+1\) mod \(k_j\). Then \(\{e^{2\pi i v X_j/k_j}, v\in \mathbb Z_{k_j}\}\) is an orthonormal basis for the law of \(X_j\), and so
\(\{\chi_u=\prod_j e^{2\pi iu_j X_j/k_j}\ :\,u\in \prod_j \mathbb Z_{k_j}=\mathcal U\}\) forms an orthonormal basis for the joint law of the \(X_j\) over the complex numbers. So we can write
$$f(X_1,\ldots, X_n)=\sum_{u\in \mathcal U}c_u\chi_u$$ for some coefficients \(c_u\). We also have
$$f(Y_1,\ldots, Y_n)=\sum_{u\in \mathcal U}\eta_uc_u\chi_u,$$
where \(\eta(u)=\prod_j \eta_j^{u_j}\), \(\eta_j=\exp(2\pi/k_j)\). We have
$$
Z_n-Z_0=\sum_{u\in \mathcal U} (1-\eta_u)c_u\chi_u.
$$
Similarly, we have
$$
Z_{j}-Z_{j-1}\stackrel{d}{=}f(X_1,\ldots,X_n)-f(X_1,\ldots,X_{j-1}, Y_{j},X_{j+1},\ldots X_n)=\sum_{u\in \mathcal U}(1-\eta_j^{u_j})c_u\chi_u.
$$
Squaring, taking expectations and using orthonormality, the desired inequality reads
$$
\sum _{u\in \mathcal U}|c_u|^2|1-\prod_j\eta_j^{u_j}|^2\le \rho(k)\sum _{u\in \mathcal U}|c_u|^2\sum_j|1-\eta_j^{u_j}|^2,
$$
which follows by summing \(|c_u|^2\) times the inequality
$$
|1-\prod_j\eta_j^{u_j}|^2\le \rho(k)\sum_j|1-\eta_j^{u_j}|^2,
$$
which translates to
$$
\sin^2(\sum_{j} \pi u_j/k_j)\le \rho(k)\sum_{j}\sin^2(\pi u_j/k_j)
$$
holding by definition of \(\rho(k)\).

For general identically distributed pairs \((X_j,Y_j)\) on \([k_j]\), we note that the joint mass function \(p_j(x,y)\) defines a sourceless flow on \([k_j]\). By the Krein-Milman theorem, the flow can be decomposed as a convex combination of unit directed flows over cycles. Since the desired inequality is linear in each of the joint laws \((X_j,Y_j)\) (but not in \(f\)), it suffices to show it for each component. Denoting the cycle lengths by \(k_j'\), and relabeling, this is exactly what was shown in the first part of the proof.
\end{proof}

\end{document}
